% TeX 格式 - 第一章 学生会线

\begin{document}

\label(chapter1_school)

\scene[student_council_room.png, peaceful_theme.mp3]

\char[sakura][smile, center]

\s 你好！欢迎来到学生会！

\char[yuki][gentle, left]

\y 啊，新同学也来了。欢迎欢迎。

\p 白雪同学也在啊。

\y 我是学生会的书记。樱井是副会长。

\s[smile] 没错！我们学生会可是学校的重要组织！

\p （学生会...应该能了解到一些学校的事情。）

\y 对了，新同学，有件事想请你帮忙。

\choice
\case[ask_for_help][+y, +t] 什么事？
\case[refuse_help][-y, -s] 我很忙...
\end-choice

% ============================================
% 分支：帮忙
% ============================================

\label(ask_for_help)

\y 最近学校里发生了一些奇怪的事情。

\s 有学生说在旧校舍附近看到了奇怪的光。

\p 奇怪的光？

\y 我们想调查一下，但是旧校舍一直锁着...

\s 如果你愿意帮忙的话，我们会很感激的！

\choice
\case[agree_to_help][+t] 我愿意帮忙
\case[worry_about_danger][+s] 这听起来很危险...
\end-choice

% ============================================
% 分支：拒绝帮忙
% ============================================

\label(refuse_help)

\s 这样啊...

\y 没关系，我们理解。

\p （他们的表情...好像很需要帮助。）

\choice
\case[ask_for_help] 等等，我改变主意了
\case[leave_student_council] 告辞了
\end-choice

% ============================================
% 同意帮忙
% ============================================

\label(agree_to_help)

\s[happy] 太好了！谢谢你！

\y 有你帮忙就放心多了。

\s 黑崎同学好像知道怎么进入旧校舍...

\p （黑崎？他确实说过知道怎么进去...）

\goto(investigation_start)

% ============================================
% 担心危险
% ============================================

\label(worry_about_danger)

\s 别担心！我们只是想看看是什么情况。

\y 如果你害怕的话，可以和我们一起行动。

\goto(investigation_start)

% ============================================
% 调查开始
% ============================================

\label(investigation_start)

\scene[school_hallway.png]

\p 好，我来帮忙调查。

\char[sakura][happy, center]

\s 太好了！那我们明天放学后在学校门口集合！

\fadeout(1.0)

\load(chapter2_investigation)

% ============================================
% 离开学生会
% ============================================

\label(leave_student_council)

\scene[school_corridor.png]

\p （也许我不该卷入这件事...）

\fadeout(1.0)

\load(chapter1_home)

% ============================================
% 额外场景：深入了解
% ============================================

\label(sakura_secret_talk)

\scene[, mystery_theme.mp3]

\s 其实...我有件事想告诉你。

\p 什么事？

\s 这所学校...隐藏着一个秘密。

\s 关于"记忆"的秘密...

\p 记忆？

\s 我不能说太多...但是请小心黑崎同学。

\s 他...不是普通人。

\p （黑崎同学？他到底是谁？）

\set(clues_found += 1)

\goto(investigation_start)

% ============================================
% 额外场景：白雪的情报
% ============================================

\label(yuki_information)

\y 我查到了一些资料...

\y 十年前，这所学校曾经进行过一项秘密实验。

\p 秘密实验？

\y 代号"星之子计划"。

\y 目的是...操纵人类的记忆。

\p 操纵记忆？！

\y 实验后来被终止了，但是...

\y 据说有些实验对象还在学校里。

\set(clues_found += 2)
\set(knows_secret = true)

\p （实验对象...还在学校里？）

\goto(investigation_start)

\end{document}
